% Full coherent reductions and conditional-input proofs; fixed original kernel.
\section{Proofs for the coherent main reductions}\label{cmr:details}
Throughout, squarefree and prime mean ideal squarefree and prime.
The original symbol is extended by zero, the complete good annulus is
unchanged, and the same complex $W$ is used for every row. The common
smooth divisor $f$ is an incidence variable after ideal M\"obius
inversion; it is not the original column-pair gcd sector $h=1$.

\subsection{Removing only the auxiliary smooth restriction}
The high-window support of a rough-prime term satisfies
\[
 N(fa),N(fb)<2D/U\le7U.
\]
The accepted exact coefficient identity~\eqref{rma:coefficient} therefore
applies to these squarefree products, including the unit:
\begin{equation}\label{cmr:short-coefficient}
 A_U(g)=\mu(g)(2\mathbf1_{Ng\le U}-1),\qquad Ng\le7U.
\end{equation}
For completeness, every proper divisor of a nonunit good $g$ has norm
at most $Ng/7\le U$. The unrestricted ordered disjoint-divisor sum is
$\prod_{p\mid g}(1-1-1)=\mu(g)$. If $Ng>U$, precisely the two
assignments $(g,1),(1,g)$ are omitted, each with weight $\mu(g)$.
If $Ng\le U$ none is omitted; $g=1$ is direct. This does not simplify
the general annular $K_0$ coefficients, whose norms can reach $2D$.

For a nonunit residual pair put $M=\max(Na,Nb)\ge7$.
On active rough-prime support, with $F=Nf$ and $P=Np>U$,
$FMP<2D\le7U^2$, hence $F<U$. Such $f$ is automatically $U$-smooth.
Thus dropping its smoothness changes only $a=b=1$.
There $F<2D/P<7U$. If a squarefree good $f$ contains a prime $q$
with $Nq>U$, any additional nonunit factor would give $F>7U$.
The omitted nonsmooth $f$ is consequently exactly $q$, with
$\mu(q)=-1$ and $|A_U(q)|=1$. This proves the sign and all masks in
\eqref{cmr:Ptwo}. The outer $p$ was already deliberately unmasked in
$K_S$; only $q\nmid u$ comes from the original $f$-row mask.
In particular $p=q$ gives a permitted auxiliary $p^2$ term, not an
original squarefree column.

A nonzero pair has both prime norms below $2D/U$.
Using~\eqref{rma:lattice}, the number of ordered pairs is at most
$(32D/U)^2$. Therefore
\[
 P_2/D\le1024B^2R_w D/U^2\le3584B^2R_w.
\]
Substitute $K_S=K_{\rm all}+P_2$ into~\eqref{rma:clean}.
The old critical error $233216B^2R_w$ and the new charged
$3584B^2R_w$ give $236800B^2R_w$ in~\eqref{cmr:auxexact}.
The original $C_2$ remains unordered and distinct, with both row
masks; it is not this ordered correction. No lower bound for the full
indefinite raw slice is inferred from either nonnegative correction.

\subsection{An exact two-jump masked Mertens functional}\label{cmr:abel}
Expand the nonunit residual part of $K_S$ before any absolute sum in
$f$, a conductor or a row. Retain
\[
 \mathcal T=S\cup\operatorname{supp}(abu),\qquad
 \eta_{a,b}(u)=\nu(a)\overline{\nu(b)}
                    \chi_a(u)\overline{\chi_b(u)}.
\]
The ideals $a,b$ are individually squarefree but may overlap.
Every common support prime remains in $\mathcal T$, even if their
quotient character loses that prime. The rough $p$ in $K_S$ may
belong to $\mathcal T$; this is its defined auxiliary enlargement.
Only the common $f$ phases cancel, leaving the full $(f,\mathcal T)=1$
mask. By~\eqref{cmr:short-coefficient}, the remaining coefficient is
\[
 \mu(f)A_U(fa)A_U(fb)
  =\mu(f)\mu(a)\mu(b)
     (2\mathbf1_{FNa\le U}-1)(2\mathbf1_{FNb\le U}-1).
\]

Put $\alpha=\min(Na,Nb)$, $M=\max(Na,Nb)$ and
\[
 X=\frac D{PM},\quad r=\frac\alpha M,\quad
 c_0=\frac U M,\quad d_0=\frac U\alpha.
\]
At $r\le1/2$ the profile product vanishes, including equality.
Otherwise let $\Phi_r(v)=W(v)\overline{W(rv)}$ when $Na=M$,
and use its conjugate in the reversed orientation. Set
$h(t)=\Phi_r(t/X)$, supported on $[X/r,2X]$ with zero endpoints.
The exact cutoff product is
\[
 (2\mathbf1_{FNa\le U}-1)(2\mathbf1_{FNb\le U}-1)
       =1-2\mathbf1_{c_0<F\le d_0}.
\]
No half-weight occurs at a norm cutoff. With the closed-prefix sum
\[
 M_{\mathcal T}(y)=\sum_{\substack{Nf\le y\\(f,\mathcal T)=1}}\mu(f),
 \qquad M_{\mathcal T}(y)=0\quad(y<1),
\]
the scalar factor, multiplied outside by
$\mu(a)\mu(b)\eta_{a,b}(u)$, is
\begin{equation}\label{cmr:scalar-kernel}
 \mathcal I_{a,b,u}(P)=
 \sum_{(f,\mathcal T)=1}\mu(f)h(Nf)
                   (1-2\mathbf1_{c_0<Nf\le d_0}).
\end{equation}
Ordinary finite Abel summation gives, for every real $y$,
\[
 \sum_{\substack{Nf\le y\\(f,\mathcal T)=1}}\mu(f)h(Nf)
   =h(y)M_{\mathcal T}(y)-\int_0^y M_{\mathcal T}(t)h'(t)\,dt.
\]
Subtracting the two closed prefixes proves the endpoint-sensitive identity
\begin{equation}\label{cmr:two-jump}
 \begin{split}
 \mathcal I_{a,b,u}(P)={}&-\int_0^\infty M_{\mathcal T}(t)h'(t)\,dt
              +2\int_{c_0}^{d_0}M_{\mathcal T}(t)h'(t)\,dt\\
       &+2h(c_0)M_{\mathcal T}(c_0)
                   -2h(d_0)M_{\mathcal T}(d_0).
 \end{split}
\end{equation}
When $c_0=d_0$ all extra terms cancel. A boundary outside the profile
has $h=0$. Combining the two integrals gives a multiplier of absolute
value one almost everywhere, rather than three separate derivative
costs. Log-coordinate differentiation yields
\[
 \int|h'(t)|\,dt\le2\log2\,M_1(W)^2,
 \qquad 2|h(c_0)|+2|h(d_0)|\le4B^2.
\]
Consequently
\begin{equation}\label{cmr:abel-variation}
 |\mathcal I_{a,b,u}(P)|
 \le6M_1(W)^2\sup_{X\le t\le2X}|M_{\mathcal T}(t)|.
\end{equation}
The full outer residual and row sum is still coherent.

A useful explicit uniformity quantity is
\[
 m_Q(Z)=\sup_{\substack{E\ \mathrm{finite\ good\ prime\ ideals}\\
                       N\prod_{p\in E}p\le Q}}
         \ \sup_{t\ge Z}\frac{|M_{S\cup E}(t)|}{t}.
\]
It always satisfies $m_Q(Z)\le16$ by~\eqref{rma:lattice}.
For the original residuals,
\[
 N\operatorname{rad}(abu)_{S^c}
 \le Na\,Nb\,Nu\le(2D/P)^2H\le14DH.
\]
Restrict the expanded nonunit $K_S$ to $X\ge Z\ge1/2$.
This cuts the residual $f$-length $X$, not $F=Nf$ in the tail below.
Equation~\eqref{cmr:abel-variation} bounds each scalar by
$12M_1(W)^2Xm_{14DH}(Z)$.
For fixed $P$ put $V=D/P$. Ideal counting gives
\[
 \#\{(a,b):\max(Na,Nb)\le t\}\le256t^2,
 \quad
 \sum_{\max(Na,Nb)\le T}\frac1{\max(Na,Nb)}\le512T.
\]
The second inequality follows by partial summation and is vacuous for
$T<1$. Since $M\le V/Z$, the absolute residual cost is at most
$6144M_1(W)^2m_{14DH}(Z)V^2/Z$.
Also $\sum_{Np>U}(Np)^{-2}\le32/U$, again by ideal counting.
Only now taking the residual, prime and row triangle proves
\begin{equation}\label{cmr:mertens-budget}
 \frac{|K_S[\mathrm{nonunit};X\ge Z]|}D
 \le196608M_1(W)^2R_w\frac D{UZ}\,m_{14DH}(Z).
\end{equation}
No arbitrary-weight sieve is hidden here: the only row input is
$|\eta_{a,b}(u)|\le1$ and $R_w$.
A \emph{hypothetical} uniform estimate
$|M_{S\cup E}(t)|\ll_{\theta,\epsilon,S}
 (N\prod_{p\in E}p)^\epsilon t^\theta$, $0<\theta\le1$,
would give, for $Z\ge1$, the restricted bound
\[
 \frac{|K_S[\mathrm{nonunit};X\ge Z]|}D
 \ll M_1(W)^2(DH)^\epsilon R_w(D/U)Z^{\theta-2}.
\]
No $\theta<1$ hypothesis is assumed.
On the region $2X\le7$, every nonunit good $f$ has zero profile
(the possible norm-$7$ endpoint is zero). Hence only $f=1$ survives,
and~\eqref{cmr:scalar-kernel} is exactly
$h(1)(1-2\mathbf1_{c_0<1\le d_0})$.
For instance $a=b$ gives $|W(1/X)|^2$.
Thus even a strong long-sum Mertens estimate cannot supply cancellation
in this one-term region after the outer residual triangle. This is a
limitation of that scalar route, not a lower bound for the complete
signed sum; other residual and row sectors may cancel it.

\subsection{Deleted-prime transfer and its fixed-field subpower input}
\label{cmr:deleted-transfer}
The deletion has an exact fixed-zeta interpretation. Put
$M_K(x)=\sum_{Na\le x}\mu(a)$, and let $\mathcal D_{\mathcal T}$ be
\emph{all} ideals supported on $\mathcal T$, including prime powers.
For $\sigma>0$ set
\[
 E_{\mathcal T}(\sigma)=
       \prod_{p\in\mathcal T}(1-(Np)^{-\sigma})^{-1}.
\]
Absolute Euler multiplication initially on $\operatorname{Re}s>1$
gives the coefficient identity and hence the finite convolution
\begin{equation}\label{cmr:deletion-convolution}
 M_{\mathcal T}(x)=
       \sum_{\substack{d\in\mathcal D_{\mathcal T}\\Nd\le x}}
                         M_K(x/Nd).
\end{equation}
Restricting $d$ to squarefree ideals would be false.
This is an artificial finite Euler deletion, not a new primitive
Hecke conductor. Its finite factors have possible poles only on
$\operatorname{Re}s=0$.

Choose $C_K$ with $A_K(x)\le C_Kx$, and define
$\rho_K(y)=\sup_{v\ge y}|M_K(v)|/(C_Kv)\le1$.
For $x\ge1$ and $0<\delta<1$, split~\eqref{cmr:deletion-convolution}
at $Nd=\sqrt x$. On the first part use $\rho_K(\sqrt x)$ and
$\sum(Nd)^{-1}\le E_{\mathcal T}(1-\delta)$.
Rankin's inequality bounds the second positive sum by
$x^{-\delta/2}E_{\mathcal T}(1-\delta)$. Thus
\begin{equation}\label{cmr:deletion-bound}
 |M_{\mathcal T}(x)|\le C_KxE_{\mathcal T}(1-\delta)
                    \{\rho_K(\sqrt x)+x^{-\delta/2}\}.
\end{equation}
Its minimum with the trivial $C_Kx$ bound is always valid.

For clarity, the fixed-field analytic input used to estimate $\rho_K$
is separate from this elementary transfer. The already pinned
\href{https://doi.org/10.2140/ant.2019.13.1039}{Thorner--Zaman,
Theorem~3.1} gives the classical Dedekind zero-free region, with a
possible single fixed exceptional real zero. Specialize only to the
trivial extension $K/K$, of absolute degree $2$ and discriminant $3$.
Shrink the fixed constant so any fixed exceptional zero is outside the
chosen region, retaining its dependence in the fixed threshold.
For $|t|\ge2$, the Euler bound just right of $1$, fixed-degree
functional equation/convexity in a strip of width $O(1/\log(|t|+3))$,
and Borel--Carath\'eodory for an analytic logarithm on smaller zero-free
disks give
\[
 |1/\zeta_K(\sigma+it)|\ll_K\log^A(|t|+3),\qquad
 \sigma\ge1-c_K/\log(|t|+3).
\]
Indeed the upper bound for $\operatorname{Re}\log\zeta_K$ and its value
at a center to the right of $1$ are $O(\log\log(|t|+3))$; a smaller
disk has the same bound for its absolute logarithm.
Bounded $t$ is handled in a fixed zero-free neighborhood of $1$,
where the reciprocal is holomorphic also at the pole of $\zeta_K$.
No logarithm is continued through that pole.

Here is the Perron error needed for the deduction. The planar
fundamental-cell count in this class-number-one field gives
$A_K(x+h)-A_K(x)\ll_K h+\sqrt x+1$ for $h\ge0$.
Apply truncated Perron to $1/\zeta_K$ at
$x'=\lfloor x\rfloor+1/2$, whose summatory value equals $M_K(x)$.
The coefficients are bounded in absolute value by the ideal counts.
Dyadic annuli around $x'$ control the near-integer terms by
$x\log^A(2x)/T+\sqrt x\log^A(2x)$; the remote terms have the first
cost. Shift to $1-c_K/\log(T+3)$, which lies inside the shrunken
zero-free region for every $|t|\le T$.
There is no residue at $1$, since the reciprocal has a zero there.
The vertical and horizontal bounds give
\[
 |M_K(x)|\ll_K x\log^A(2x)
 \left\{\exp\left(-\frac{c_K\log x}{\log(T+3)}\right)+T^{-1}\right\}
       +\sqrt x\log^A(2x).
\]
Choosing $T=\exp(c'_K\sqrt{\log(2x)})$ and absorbing logarithms proves
\begin{equation}\label{cmr:fixed-mertens}
 |M_K(x)|\ll_K x\exp(-c''_K\sqrt{\log(2x)}).
\end{equation}
Small $x$ is absorbed in the fixed constant. This is fixed-field
subpower cancellation; it asserts no fixed power or uniformity over
varying primitive conductors or fields.

The remaining deletion factor can be quantified elementarily. Put
$R_{\mathcal T}=N\prod_{p\in\mathcal T\setminus S}p$ and
$Z_0=\log(e^3R_{\mathcal T})\ge3$. For $0<\delta\le1/2$, split the
prime sum for $\log E_{\mathcal T}(1-\delta)$ at $Z_0^2$.
Ideal counting bounds the small part by
$O_K(Z_0^{2\delta}\log Z_0)$. For the large part use
$\sum_{p\in\mathcal T\setminus S}\log Np\le\log R_{\mathcal T}$
and the monotonicity of $v^{1-\delta}\log v$; its bound is
$O(\log R_{\mathcal T}/(Z_0^{2-2\delta}\log Z_0))=O(1)$.
Fixed $S$ factors cost a fixed constant. Hence
\[
 \log E_{\mathcal T}(1-\delta)
       \le C_{K,S}\{1+Z_0^{2\delta}\log Z_0\}.
\]
If $\delta\log Z_0\le1$, this is $E_{\mathcal T}(1-\delta)\ll Z_0^A$.
Taking $\delta=\min\{c/\sqrt{\log(2x)},1/\log Z_0,1/2\}$ in
\eqref{cmr:deletion-bound} proves, for $x\ge3$,
\begin{equation}\label{cmr:uniform-mertens}
 |M_{\mathcal T}(x)|\ll_{K,S}xZ_0^A
 \exp\left(-c\min\left\{\sqrt{\log(2x)},
                          \frac{\log(2x)}{\log Z_0}\right\}\right).
\end{equation}
For $R_{\mathcal T}\le D^C$ and $x\ge D^v$, fixed $C,v>0$,
this gives a uniform $x\log^A(2D)\exp(-c\sqrt{\log(2x)})$ bound.
It is useful on long scalar ranges, but gives no power saving and
no cancellation in the unit-$f$ range above. The exact product bound
\eqref{cmr:deletion-bound} retains every raw deletion before any
convenient larger estimate.

\subsection{The joint matrix, raw zeros and its compression}
\label{cmr:joint-proof}
All spaces here carry ordinary counting $\ell^2$ norms. On the fixed
finite short-column space of Proposition~\ref{cmr:joint},
$g=fa$ is squarefree precisely when $(f,a)=1$.
For $(f,u)=1$, multiplicativity gives
\[
 (\mathsf R\mathsf P_fb)(u)
   =\sqrt{w_u}\nu(f)\chi_f(u)T_f(t;u).
\]
If $(f,u)>1$, every $g$ in that projected column support has the
original zero $\chi_g(u)=0$. Thus its row is exactly zero, rather
than being restored through a primitive character. Squaring proves
$G_f(t)=\|\mathsf R\mathsf P_fb(t)\|^2$ and
\eqref{cmr:joint-form}. The Gram entries remain
\[
 (\mathsf R^*\mathsf R)_{g,h}
      =\sum_u w_u\overline{\chi_g(u)}\chi_h(u).
\]
No coherent orientation mean is removed.

The divisor sum in each matrix entry is the exact Schur factor
\[
 C_U(g,h)=\sum_{\substack{f\mid(g,h)\\f\ U\text{-smooth}}}\mu(f)
       =\mathbf1_{(g,h)\text{ has no prime norm }\le U}.
\]
On active $t\ge U$ support, $Ng,Nh<7U$. Any such $g$ containing a
rough prime must be that prime alone. It follows that
\begin{equation}\label{cmr:schur}
 C_U(g,h)=\mathbf1_{(g,h)=1}
       +\mathbf1_{g=h=p,\ p\notin S,\ Np>U}.
\end{equation}
For the projection $\Pi_{\mathcal P}$ onto rough-prime columns,
$\mathsf P_f\Pi_{\mathcal P}=0$ for every nonunit smooth $f$ and
$\mathsf P_1\Pi_{\mathcal P}=\Pi_{\mathcal P}$. Therefore
\begin{equation}\label{cmr:compression}
 \Pi_{\mathcal P}\mathcal M_U\Pi_{\mathcal P}
     =\Pi_{\mathcal P}\mathsf R^*\mathsf R\Pi_{\mathcal P}.
\end{equation}
This is a quadratic-form \emph{compression}, not an invariant
restriction: in general $\mathcal M_U\Pi_{\mathcal P}$ and
$\Pi_{\mathcal P}\mathcal M_U$ differ. Off-block couplings and other
columns can cancel its contribution. It is no lower bound for full
$Q$, the prime error or the original raw slice.

On any common active length dyad $[L,2L]$, a true squared row norm
$\lambda_*(L)$ gives, for arbitrary coefficients $c$ there,
\[
 |c^*\mathcal M_Uc|
 \le\sum_f\|\mathsf R\mathsf P_fc\|^2
 \le\lambda_*(L)\sum_g d_U(g)|c_g|^2,
 \qquad d_U(g)=\#\{f\mid g:f\ U\text{-smooth}\}.
\]
This is a positive incidence-frame majorant, not M\"obius cancellation.
Since $g$ is squarefree, $d_U(g)\le2^{\omega(g)}\ll_\eta(Ng)^\eta$.
Self-adjointness bounds the compressed operator norm by the same
$\lambda_*(L)\max d_U$ and permits polarization for cross terms.
For $L=D/t$, the exact short coefficient and ideal counting give
\[
 \|b(t)\|^2\ll B^2L,\qquad
 \|tb'(t)\|^2\ll M_1(W)^2L.
\]
The direct unit endpoint covers $1/2\le L<1$; no column survives below
$1/2$. Differentiating the fixed matrix, not a moving support matrix,
gives $Q'(t)=2\operatorname{Re}(b'(t)^*\mathcal M_Ub(t))$ and
\begin{equation}\label{cmr:variation}
 |Q(t)|+|tQ'(t)|\ll_\eta(HD)^\eta
                     M_1(W)^2\lambda_*(D/t)D/t.
\end{equation}
Only here may the classical specialization be inserted, in exactly
its original natural-row scope. For arbitrary weights a true
weighted norm is needed; the elementary Frobenius alternative is
$\lambda_*(L)\ll R_wL$.

\subsection{Exact Stieltjes boundary and squarefunction criterion}
\label{cmr:stieltjes-proof}
The finite $Q$ vanishes for $t\ge2D$, including its endpoint.
Stieltjes integration over $(U,\infty)$ counts precisely $Np>U$;
comparing its prime measure with $dt/\log t$ proves
\eqref{cmr:primeerror}. Changing the normalization of $\operatorname{Li}$
by a constant changes the two boundary/integral terms by opposite
amounts and changes no identity. We fix $\operatorname{Li}_2$ here.
The original raw slice is consequently
\[
 S_1^{AA}=\frac{K_0-\mathcal J}D+\Delta+\frac{C_2}D
                                   -\frac{\mathcal E_{\rm PIT}}D.
\]
It keeps the original masks and the full mean. The prime measure
contains only fixed $S$, never a growing deleted set
$S\cup\operatorname{supp}(abu)$.

Let $\delta_S(U)=\sup_{t\ge U}|E_S(t)|/t$.
Use~\eqref{cmr:variation} and substitute $L=D/t$ to obtain
\begin{equation}\label{cmr:pointwise-budget}
 \frac{|\mathcal E_{\rm PIT}|}D
 \ll_\epsilon(HD)^\epsilon M_1(W)^2\delta_S(U)
 \left\{\lambda_*(D/U)+
        \int_{1/2}^{D/U}\lambda_*(L)\,\frac{dL}{L}\right\}.
\end{equation}
One may keep $\max d_U$ explicitly instead of absorbing it into a
subpower; no uniform logarithmic improvement is thereby claimed.
With the current physical input of Proposition~\ref{phy:short}, the minimum of
its three branches remains inside this integral. All constants and
internal small exponents are chosen uniformly on the neighboring
capacity-feasible blocks.

For $I_T=[T,2T]\cap[U,2D]$, let $\lambda^\#(D/T)$ dominate
$\lambda_*(L)$ on $L\in[D/(2T),D/T]$, using the direct endpoint
when necessary. Equation~\eqref{cmr:variation} implies
\[
 \|Q'\|_{L^2(I_T)}\ll(HD)^\eta M_1(W)^2
                         \lambda^\#(D/T)D/T^{3/2}.
\]
Our weighted squarefunction convention gives
$\|E_S\|_{L^2(I_T)}\le(2T)^{3/2}\mathfrak d_2(T)$.
Cauchy on each dyad and the retained endpoint prove
\eqref{cmr:square-budget}. The weighted definition is comparable by
absolute constants to $T^{-3/2}\|E_S\|_{L^2([T,2T])}$.
Global Cauchy across unrelated dyads need not improve this bound.
The sufficient target condition is that its braces be
$\ll H(HD)^\epsilon$, including the pointwise term. It is an explicit
input to investigate, not an established estimate or an independence
assumption between $E_S$ and the signed derivative.

\subsection{Mellin Parseval retains equal-norm coherence}\label{cmr:mellin}
Extend only the same finite short-column vector to $t>0$ for this
identity; below $U$ it need not equal the unrestricted original
$T_f$ with additional longer columns. No signed equality for $K_S$
is enlarged below $U$. For fixed $f,u$ let
$c_g=A_U(g)\nu(g)\chi_g(u)\mathbf1_{f\mid g}$, and use
\[
 \widehat W(i\tau)=\int_0^\infty W(x)x^{i\tau}\,\frac{dx}{x}.
\]
Ordinary Mellin Plancherel for the finite amplitude gives exactly
\begin{equation}\label{cmr:parseval}
 \begin{split}
 &\int_0^\infty\left|\sum_g c_gW(Ng\,t/D)\right|^2\frac{dt}{t}\\
 &=\frac1{2\pi}\int_{\mathbb R}|\widehat W(i\tau)|^2
                   \left|\sum_gc_g(D/Ng)^{i\tau}\right|^2d\tau.
 \end{split}
\end{equation}
For the derivative of this \emph{amplitude}, replace
$|\widehat W|^2$ by $\tau^2|\widehat W|^2$.
This does not assert a diagonal Mellin formula for the square of the
quadratic derivative $Q'$.
Finite signed summation in $f$ and the actual rows preserves the
matrix form at the norm-twisted coefficient vector. No absolute-value
or centering operation is inserted into that equality.

If distinct ideals $g,h$ have $Ng=Nh$, their radial profiles and all
their $t$ derivatives are identical. Their radial cross correlation
is the same nonzero profile norm as the diagonal when $W\ne0$.
The Mellin phases $(Ng)^{-i\tau}$ also coincide; distinct conjugate
split ideals provide actual such orientations. Thus neither radial
Parseval nor the derivative removes their coherent mean.
The safe decomposition is
\[
 \sum_r a_r\Gamma_r
   =\Gamma_{\rm av}\sum_ra_r
                 +\sum_ra_r(\Gamma_r-\Gamma_{\rm av}).
\]
Both terms remain. The already recorded fixed-profile example in
Section~\ref{sw:core} has centered variance sum $6$, coefficient norm
squared $2$, but full scalar squared $36>12$; it is not rerun or
identified with this prime error. It illustrates why a variance-only
replacement is invalid. Applying an existing arbitrary-coefficient
norm at every $\tau$ in~\eqref{cmr:parseval} leaves its operator
power unchanged; the twist has modulus one.

\subsection{Proof of the inherited smooth-divisor localization}
\label{cmr:tail-proof}
Now specialize only to the original unweighted $A_6=1$ two-pivot
rows. The exact $K_0$ positive summand is
\[
 G_f(1)=\sum_{\substack{u\in\mathcal U\\(f,u)=1}}
             |T_f(1;u)|^2.
\]
Here $A_U(fa)$ is in the general annular domain and stays unfactored.
The divisor-assignment bound and ideal counting give its coefficient
squared norm $\ll_\eta B^2D^\eta L$, $L=D/F$.
The accepted powered subsequence norm~\eqref{sw:divisor-budget} gives
\[
 G_f(1)\ll(HD)^\eta B^2
 \left\{LR+\frac{L^2R}{M_{\rm row}}
              +\frac{L^{5/3}R}{M_{\rm row}^{1/3}}\right\}.
\]
For $1/2<L<1$ count the unit residual directly. Its
$O(RD^\eta B^2)$ cost is covered by the same expression up to a
constant; for $L=1/2$ the endpoint is zero. No subunit sieve is used.
Now take the triangle only in the tail $F>Y$, and enlarge the
nonnegative squarefree smooth-ideal upper sums to all ideals.
Partial summation from ideal counting yields
\[
 \sum_{Y<F<2D}F^{-1}\ll\log(2D),\quad
 \sum_{F>Y}F^{-2}\ll Y^{-1},\quad
 \sum_{F>Y}F^{-5/3}\ll Y^{-2/3}.
\]
Division by $D$ and absorption of logarithms and internal assignment
powers prove~\eqref{cmr:tailbound}. Every rough-prime $K_S$ term has
$FPNa>FU$, so $F>Y\ge2D/U$ forces zero profile even in its auxiliary
nonsquarefree cases. The $C_2$ and $P_2$ unit corrections are not
silently removed by this tail split.

When $R\ll H$, solving the last two monomial inequalities for a
sufficient $Y$ gives the displayed cutoff in Proposition~\ref{cmr:tail}.
At the mixed witness it is enough to take $Y=D^{13/24}$ with fixed
capacity slack, giving tail exponents $7/12,5/8,3/4$.
At the pure witness $Y=D^{5/8}$ gives $3/4,3/8,3/4$.
Both exceed the rough support exponent $1/2$ for large $D$.
These are inherited sufficient upper cutoffs, not energy lower bounds
or assertions of current optimality after the new physical input.
All small smooth divisors, including $f=1$, remain; no signed main
estimate follows from the positive tail. The current whole-form
comparison in Proposition~\ref{phy:signed} is a different consequence of the
new external norm input.

\subsection{Proof of the conditional zero-free input diagnostic}
\label{cmr:strip-proof}
Assume exactly the all-scale, all-$\epsilon$ hypothesis of
Theorem~\ref{cmr:conditional-strip}. For $\operatorname{Re}s>1$ put
\[
 F_S(s)=\int_2^\infty E_S(t)t^{-s-1}\,dt.
\]
On a dyad, Cauchy gives
\[
 \int_T^{2T}|E_S(t)|t^{-\sigma-1}\,dt
 \le\mathfrak d_2(T)
       \left(\int_T^{2T}t^{1-2\sigma}\,dt\right)^{1/2}
 \le C_\sigma\mathfrak d_2(T)T^{1-\sigma}.
\]
At $\sigma=1$ use the logarithmic limit of the finite integral.
On a compact subset of $\sigma>1-\alpha$, choose $\epsilon$ smaller
than its distance to that boundary. The dyadic majorant is a
summable geometric series, also after any fixed power of $\log t$.
The same integral therefore converges normally, with all holomorphic
derivatives, on $\sigma>1-\alpha$. Its bounded initial interval is
entire. No contour has crossed an unknown zero.

For $\sigma>1$, the prime-ideal series satisfies
\[
 P_S(s)=\sum_{p\notin S}(Np)^{-s}
       =L_0(s)+sF_S(s),\qquad
 L_0(s)=\int_2^\infty\frac{t^{-s}}{\log t}\,dt.
\]
In this field $\pi_K^S(2)=0$ and $\operatorname{Li}_2(2)=0$, so no
lower boundary atom is missing. Direct differentiation gives
$L_0'(s)=-2^{1-s}/(s-1)$. Only this meromorphic derivative is
continued; no single-valued logarithmic primitive is asserted through
$s=1$. Consequently
\[
 -P_S'(s)=\frac{2^{1-s}}{s-1}-(sF_S(s))'.
\]
For $\zeta_K^S(s)=\zeta_K(s)\prod_{p\in S}(1-(Np)^{-s})$, its
initial logarithmic derivative is
\[
 -\frac{(\zeta_K^S)'(s)}{\zeta_K^S(s)}
   =-P_S'(s)+H_S(s),\quad
 H_S(s)=\sum_{p\notin S}\sum_{k\ge2}
                          \log(Np)(Np)^{-ks}.
\]
The higher-prime-power series is normally holomorphic for
$\sigma>1/2$: its geometric tail is bounded on compact subsets by
$C\sum_p\log(Np)(Np)^{-2\sigma}$, an absolutely convergent ideal
series. Inert prime ideals are included, not mistaken for absent
rational primes. The right side has no possible pole in
$\sigma>\max\{1-\alpha,1/2\}$ except $s=1$.
Standard continuation of the fixed Dedekind zeta, or
$\zeta_K(s)=\zeta(s)L(s,\chi_{-3})$, identifies it with the actual
logarithmic derivative. A zero elsewhere would give a pole with
nonzero integer residue, impossible there. Finite deleted Euler
factors vanish only on $\sigma=0$, so they change no zero in this
region. This proves Theorem~\ref{cmr:conditional-strip}.
A fixed exceptional real zero, if inside the strict region, would be
excluded by the hypothetical input itself; it is not assumed absent.

For orientation, the old and current sufficient scalar power accounting
is recorded below. Every input and ensuing strip in the table is
conditional, and endpoint control is additional to a mean-square norm.
Here $\sigma=\operatorname{Re}s$ and the input column means $U^{-\alpha}$.
\begin{center}
\begin{tabular}{lccc}
certificate & budget & input $\alpha$ & strict strip\\
\hline
old pure (historical) & $5/6$ & $1/6$ & $\sigma>5/6$\\
current pure & $19/24$ & $1/12$ & $\sigma>11/12$\\
current mixed & $7/9$ & $1/18$ & $\sigma>17/18$
\end{tabular}
\end{center}
The current values use Proposition~\ref{phy:short}, not the historical pure
$5/6$ operator budget. Since $U\asymp D^{1/2}$, subtracting the
target $3/4$ from the respective PIT powers and doubling gives
$1/12$ and $1/18$. A pointwise all-scale power for $\delta_S$ would
imply the corresponding squarefunction power; the converse and its
pointwise endpoint are not supplied. The accepted fixed-field
subpower PIT is not one of these independent powers.

The new physical input does prove its stated global $9/8$ and
$A_6=1$ signed/combined $13/12$ comparisons. The coherent reductions
above do not improve them to the target, nor establish desired
positivity, general Artin holomorphy, determinant/zero-location
bridges, a varying-profile fixed-seminorm theorem or new Lean results.
Finite formal controls and source-bound reviews check bookkeeping;
they do not replace any infinite analytic hypothesis. Full primary
reading caches remain local: only public citation/version/hash
metadata accompanies these proofs.
